%% body_related.tex
%% Shared manuscript source, extracted verbatim from the historical
%% entry point paper/main.tex (lines 2998-3133) so that it and the TMLR
%% entry point submissions/tmlr_2026/main.tex read one copy of this
%% material instead of two. Content is unchanged by the extraction
%% itself. The revision the lines refer to is recorded in
%% paper/EXPOSITION_REVISION.md, which is not part of any submission
%% upload.
\section{Related Work}\label{sec:related}

The closest methods differ primarily in where they freeze features and how they
represent uncertainty.  That distinction matters here because a covariance
built from stored features answers a different question from one that replays
every observation at the current parameter.

\paragraph{Positive-cone metric geometry.}
\citet{thompson1963certain} introduced the part metric now commonly called the
Thompson metric on ordered cones.  \citet{nussbaum1994finsler} developed its
Finsler structure and applications.  On the positive-definite cone, the
normalized operator norm
$\|V^{-1/2}\dot V V^{-1/2}\|_{\rm op}$ is the corresponding tangent norm.
Our contribution is not this geometry or its path-to-endpoint comparison.  We
instantiate it as an operational confidence-transport device linking frozen
bandit confidence, current replayed GGN/Fisher curvature, approximate operators,
and certified solver widths.

\paragraph{Neural contextual bandits.}
LinUCB \citep{abbasi2011improved} and Thompson sampling
\citep{thompson1933likelihood,agrawal2013thompson} provide the statistical
backbone for many neural extensions.  In their stated algorithms, NeuralUCB
\citep{zhou2020neural} and NeuralTS \citep{zhang2021neural} add a gradient
feature to the covariance when its observation arrives; they do not replay the
entire history at each new parameter.  \citet{kassraie2022neural} separate a
kernelized NTK reference from a trained-network mean with initialization-time
uncertainty features, and prove sublinear regret for the SupNN-UCB variant under
NTK-RKHS, sufficient-exploration, and network-width assumptions.
\citet{salgia2023provably}
controls adaptivity through grouped histories and a finite-network-to-NTK
confidence argument.  \citet{deb2024contextual} take a different route through
online neural regression and contextual-bandit reductions.  These results do
not lack proofs because their covariance differs from ours; they answer
different questions.  We analyze the additional case in which every historical
Jacobian is replayed at the current parameter and the resulting operator is
solved approximately.

The NeuralLinear method maintains Bayesian linear uncertainty in learned
final-layer features \citep{riquelme2018deep}; the methods of
\citet{xu2022shallow} and \citet{bui2025variance} instead use UCB-style linear
exploration over learned deep representations.  Recomputing those last-layer
representations is not the same as transporting confidence to a full-parameter
GGN.  Our supplied-subspace result recovers a projected linear geometry when
such a subspace is fixed in advance, but it does not turn the general theorem
into a neural-linear analysis.

\paragraph{Online Bayesian and EKF-style neural bandits.}
\citet{duranmartin2022efficient} apply extended Kalman filter (EKF)
updates in a learned or random low-dimensional affine parameter subspace.
They maintain a dense or diagonal covariance in that subspace, inducing a
low-rank covariance over parameters spanning all network layers, with memory
independent of the number of rounds.  This differs from applying exact
curvature--vector products in the full parameter space.
More directly, \citet{chang2023lowrank} maintain a diagonal-plus-low-rank
approximation to the full-parameter posterior precision through recursive EKF
and online-SVD updates, and use it for Thompson sampling in an MNIST contextual
bandit.  Their LO-FI method is a one-pass compressed-posterior approach with
fixed-rank cost linear in parameter count and no replay of past data.  CC-UCB
instead relinearizes the full replay history at the current parameter and
accesses the uncompressed GGN through residual-checked CG.  Their bandit study
is empirical; our conditional UCB analysis explicitly accounts for CG
truncation, but does not establish a practical advantage over LO-FI.

\paragraph{Linearized networks and Laplace approximations.}
\citet{maddox2021fast} linearize once around a trained network and use implicit
Jacobian/Fisher products with CG for predictive inference.  Their construction
is a computational antecedent for matrix-free widths, but the reference
linearization is fixed rather than replayed online.  The Laplace approximation
places a local Gaussian posterior around a MAP solution with covariance defined
by inverse curvature \citep{mackay1992practical}.  Last-layer Laplace fixes the
backbone, omitting backbone-weight uncertainty by construction
\citep{kristiadi2020being}.  \citet{daxberger2021laplace} systematize all-weight,
subnetwork, and last-layer variants with full, low-rank, Kronecker-factored, or
diagonal curvature; structured all-layer variants use Kronecker-factored
curvature \citep{ritter2018scalable,martens2015optimizing}.  Our concern is
instead the sequential validity of a UCB score when the current relinearized
operator differs from the predictable confidence Gram.

\paragraph{Sampling-based full-parameter exploration.}
Langevin Monte Carlo Thompson sampling (LMC-TS) \citep{xu2022langevin}
explores with approximate samples of a loss-defined parameter posterior,
obtained by noisy gradient steps on the regularized loss; like CC-UCB it never
materializes a covariance matrix.  Its regret guarantee is for linear
contextual bandits; its neural-bandit evidence is empirical.
The two routes pay for exploration differently: LMC-TS through the
mixing of a per-round sampling chain, CC-UCB through CG truncation and
spectral distortion of an explicit curvature metric, which yields
deterministic widths, an optimism-style analysis, and a fixed curvature operator
reused across separate per-action CG solves.
\paragraph{Scope of efficient nonlinear exploration.}
The methods above obtain efficiency through different statistical objects:
fixed NTK features, grouped data, regression reductions, or learned last-layer
representations.  CC-UCB targets a practitioner who chooses to score with a
specific current full-parameter GGN/Fisher operator.  Its conditional theorem
exposes what that choice requires: transport from collection-time confidence,
operator distortion, and solver error.  It does not imply that this route is
computationally preferable to the alternatives.

\paragraph{Spectral sketching for linear bandits.}
\citet{kuzborskij2019efficient} analyze linear bandits where the Gram
matrix is replaced by a Frequent Directions sketch, deriving regret
bounds that depend on the spectral approximation quality.
\citet{wen2026revisiting} show that fixed-size sketch regret bounds can
degenerate to the trivial linear rate under heavy spectral tails, and propose
an adaptive dyadic block sketch.
Our spectral-distortion theorem (\Cref{thm:spectral-distortion}) instantiates the
same conceptual lever in the matrix-free GGN/CVP setting by exposing transfer
and realized-width terms along the realized trajectory.
We do not claim novelty for the sandwich-to-regret reduction itself;
the contribution is adapting it to the CG/GGN context with explicit
CG truncation and feature-drift factors.

\paragraph{GGN--vector products, Hessian spectral analysis, and influence
functions.}
Pearlmutter's algorithm \citep{pearlmutter1994fast} computes exact
Hessian--vector products; the analogous GGN--vector product has the same
Jacobian-transpose-times-Jacobian-times-vector structure
\citep{schraudolph2002fast}.
Recent work on Krylov methods for Hessian spectral analysis
\citep{granziol2026hessian} demonstrates related distributed HVP/Krylov
infrastructure at $>10^8$-parameter scale, but does not benchmark the exact
per-example GGN operator or end-to-end CC-UCB latency.
Influence functions
\citep{koh2017understanding} use inverse-curvature--vector products to
attribute predictions to training points; the CG-based implicit solve
we employ is the same computational primitive.

\paragraph{Optimal experimental design.}
Information-theoretic exploration objectives---such as maximizing
$\log\det(\bCbar_t)$ or minimizing the trace of~$\bCbar_t^{-1}$---connect
bandit exploration to the classical D-optimal and A-optimal design
criteria \citep{valko2013finite}.  CC-UCB's use of the full GGN/Fisher curvature
enables stochastic estimates of these quantities via Hutchinson-type trace
estimators \citep{ubaru2017fast}, without forming the matrix.
